\subsection{An arithmetic counterexample to eventual equality in A.2}

\begin{lemma}[The elementary two-group calculation]
\label{C:A2-wreath-calculation}
Let $V$ be a nonzero finite-dimensional vector space over $\mathbf F_2$.
Put $G=(V\oplus V)\rtimes\langle\tau\rangle$, where $\tau$ interchanges
the two factors, put $H=V\oplus V$, and define
\[
 \Sigma:G\longrightarrow V,\qquad (x,y;\epsilon)\longmapsto x+y.
\]
Then $A=\ker\Sigma$ is an elementary abelian two-group. Writing its
elements as $(b,u)\leftrightarrow(u,u;b)$ and using the section
$s(v)=(v,0;0)$ of $\Sigma$, the conjugation action is
\[
                 v\cdot(b,u)=(b,u+bv).
\]
There is a surjective $V$-module homomorphism
\[
 \phi:\mathbf F_2[V]\longrightarrow A,
 \qquad e_g\longmapsto(1,g),
\]
and hence a surjection $\mathbf F_2[V]\rtimes V\to G$ inducing the
identity on the quotient $V$. Each $\phi(e_g)$ is a conjugate of
$\tau$, and its cyclic subgroup intersects $H$ trivially.
The transfer from $G$ to its index-two abelian subgroup $H$ is
\[
             \operatorname{Ver}_{G,H}(x,y;\epsilon)
                    =(x+y,x+y).
\]
\end{lemma}
\begin{proof}
Multiplication on $A$ is componentwise addition because the two entries
of $(u,u)$ agree. Conjugation by $s(v)$ sends $(u,u;b)$ to
$(u+bv,u+bv;b)$, proving the action formula. Translation by $v$ sends
$e_g$ to $e_{g+v}$, which proves equivariance. The image contains
$(1,0)$ and $(1,g)+(1,0)=(0,g)$ for every $g$, so it is all of $A$.
The conjugates of $\tau$ by $(g,0;0)$ are $(g,g;1)$, precisely the
stated images. Their nonidentity elements have nonzero image in
$G/H\simeq C_2$.

For the transfer, use coset representatives $1,\tau$. An element
$(x,y;0)$ transfers to the product with its $\tau$-conjugate,
namely $(x+y,x+y)$. An element $(x,y;1)$ exchanges the two cosets;
its transfer is its square, again $(x+y,x+y)$. These calculations
also prove that in $A\rtimes V$ an element $(b,u;v)$ has square
$(0,bv;0)$.
\end{proof}

\begin{lemma}[A quadratic step preserving ideal classes modulo two]
\label{C:A2-quadratic-step}
Let $F$ be a totally imaginary number field, and let $S$ be a finite
set of its finite primes containing every prime above $2$. Write
\[
 C_F=\operatorname{Pic}(\mathcal O_F[1/2])/2,
 \qquad r=\dim_{\mathbf F_2} C_F>0.
\]
There is a quadratic extension $L/F$ with the following properties:
\begin{enumerate}
 \item Every prime in $S$ is unramified and inert in $L/F$.
 \item Extension of invertible ideals induces an injection
       $C_F\hookrightarrow C_L$.
 \item $\dim_{\mathbf F_2}C_L\ge 2r$.
\end{enumerate}
More precisely, we can construct a Galois extension $N/F$ containing
$L$ such that
\[
 \operatorname{Gal}(N/F)\simeq (V\oplus V)\rtimes C_2,
 \qquad \operatorname{Gal}(N/L)=V\oplus V,
 \qquad V\simeq C_F,
\]
and $N/L$ is unramified at every finite prime and splits completely
at every prime above $2$.
\end{lemma}
\begin{proof}
Let $M/F$ be the maximal elementary abelian two-extension unramified
at all finite primes and split at the primes above $2$. Global class
field theory identifies $V=\operatorname{Gal}(M/F)$ with $C_F$ by
the Artin map. We use the ideal-class existence and reciprocity
theorems in \cite[Chapter V, Theorems 3.5--3.6, p.~158]{C_MilneCFT2020},
and take the quotient in which the classes of the dyadic primes
vanish. Thus $M/F$ is everywhere unramified and $2$-split.

For each $\mathfrak q\in S$, choose one prime
$\mathfrak Q_{\mathfrak q}$ of $M$ above it. At that prime choose
a unit squareclass defining the unramified quadratic extension of
the local field $M_{\mathfrak Q_{\mathfrak q}}$; at all the other
primes of $M$ above $\mathfrak q$, choose the square class.
These instructions include every dyadic prime. Such a unit class
exists also in residue characteristic two: lift an irreducible
quadratic $T^2-T+\overline c$ from the residue field and use its
unit discriminant $1-4c$. Its roots generate the unramified quadratic
extension. The prescribed local squareclasses are open, since
sufficiently deep principal units are squares by Hensel's lemma.
Weak approximation therefore gives $b\in M^\times$ that is a unit
at all these primes and has all the prescribed squareclasses.

Choose a $V$-invariant integral ray modulus $\mathfrak m$, supported
on the primes above $S$, deep enough that multiplication by an
element congruent to $1$ modulo $\mathfrak m$ changes none of these
local squareclasses. Let $E/M$ be its ray class field. Invariance
of the modulus makes $E/F$ Galois: every $F$-automorphism of $M$
preserves the defining ray subgroup and hence its class field.
Let $g\in\operatorname{Gal}(E/M)$ be the Artin image of the ray
class of the principal fractional ideal $(b)$.

Chebotarev applied to $E/F$ supplies a prime
$\mathfrak p\notin S$, unramified in $E$, and a prime of $E$ above
it whose Frobenius is $g$; first choose its conjugacy class and then
choose the prime to obtain this representative. Since $g$ acts
trivially on $M$, the prime $\mathfrak p$ splits completely in $M/F$.
For the resulting prime $\mathfrak P$ of $M$, its ray class is
$[(b)]$. By the definition of ray equivalence there is therefore
an element $\eta\in\mathcal O_M$ such that
\[
             \mathfrak P=(\eta),\qquad \eta/b\equiv1\pmod{\mathfrak m}.
\]
Here the equality of ideals makes $\eta$ integral, and the congruence
is only at the modulus primes. This application uses
\cite[Chapter V, Theorem 3.23, p.~164]{C_MilneCFT2020}; it requires
no simultaneous inverse-Galois or embedding theorem.

Now put
\[
       \widetilde N=M\bigl(\sqrt{\sigma\eta}:\sigma\in V\bigr).
\]
The principal ideals $(\sigma\eta)=\sigma\mathfrak P$ are distinct.
Their valuations show that the classes of all $\sigma\eta$ are
linearly independent in $M^\times/M^{\times2}$. Consequently
$\operatorname{Gal}(\widetilde N/M)=\mathbf F_2[V]$, with $e_\sigma$
changing the sign of $\sqrt{\sigma\eta}$ alone. The natural action
of $V$ lifts by the formulas
\[
 g\bigl(\sqrt{\sigma\eta}\bigr)=\sqrt{g\sigma\eta}.
\]
These formulas define an action, not merely separate lifts: the
algebra with these square-root generators is a field by the
independence just proved, and permutation of its labeled generators
respects every defining relation and the multiplication law of $V$.
Thus $\operatorname{Gal}(\widetilde N/F)=\mathbf F_2[V]\rtimes V$.

Use Lemma~\ref{C:A2-wreath-calculation} to quotient by $\ker\phi$.
The resulting $N/F$ has group $G=V\wr C_2$, and the old extension
$M/F$ is the quotient $\Sigma:G\to V$. Let $L=N^H$, where
$H=V\oplus V$; then $[L:F]=2$.

We first check ramification at every prime. At a nondyadic prime
of $M$ outside the primes over $\mathfrak p$, all radicands are
units, so the local quadratic extensions are split or unramified.
At dyadic primes they are split or unramified by the imposed local
squareclasses, which are permuted under $V$. At the prime
$\sigma\mathfrak P$ the only radicand of odd valuation is
$\sigma\eta$. Thus the inertia group there in
$\widetilde N/M$ is exactly $\langle e_\sigma\rangle$; the other
unit radicands contribute only unramified extensions. Because
$\mathfrak p$ splits in $M/F$, this is also the inertia in
$\widetilde N/F$. In $G$ its image is the complement conjugate
$\langle\phi(e_\sigma)\rangle$, whose intersection with $H$ is
trivial. It follows that $N/L$ is unramified at every finite prime.
There are no real places.

We next verify inertness at each old prime and splitting above two.
Fix $\mathfrak q\in S$ and let $v\in V$ be its Frobenius in $M/F$.
If $v=0$, all the primes above $\mathfrak q$ have degree one. At
one such prime exactly one conjugate radicand has the nonsquare
unramified local class, so its Frobenius in the regular module is
$e_g$ for some $g$. Its image $(1,g)\in A$ has order two and
nonzero image in $G/H$. Thus $L/F$ is unramified inert at
$\mathfrak q$, and $N/L$ is split at the prime over it.
In particular this holds at every dyadic prime, since $M/F$ is
$2$-split.

If $v\ne0$, the decomposition subgroup in $M/F$ is
$\langle v\rangle$, of order two. At a chosen prime above
$\mathfrak q$, precisely two conjugate radicands have the
nonsquare unramified local class, with indices $g$ and $g+v$.
Therefore the Frobenius of $\widetilde N/M$ is $e_g+e_{g+v}$,
and its image under $\phi$ is $(0,v)\in A$.
If $z\in G$ is the Frobenius over $\mathfrak q$, then
$\Sigma(z)=v$ and $z^2=(0,v)$ in $A$. The square formula in
Lemma~\ref{C:A2-wreath-calculation} says that a lift $(b,u;v)$
has square $(0,bv)$; since $v\ne0$, we have $b=1$.
Hence $z$ again has nonzero image in $G/H$, proving unramified
inertness in $L/F$.

Finally, Artin reciprocity gives a surjection
$C_L\twoheadrightarrow H$ because $N/L$ is an unramified,
$2$-split, elementary abelian two-extension. Extension of ideals
corresponds to transfer; the exact compatibility is
\cite[Chapter V, Proposition 3.18 and the paragraph following it,
p.~162]{C_MilneCFT2020}. The transfer calculation gives the
commutative diagram
\[
 \begin{array}{ccc}
 C_F&\longrightarrow&C_L\\
 \big\downarrow\wr&&\big\downarrow\\
 V&\xrightarrow{\ v\mapsto(v,v)\ }&V\oplus V.
 \end{array}
\]
One may compute this diagram on ideals prime to the finite
ramification set of $N/F$; every ideal class has such a representative.
The bottom map is injective, so the top map is injective. The
surjection on the right gives $\dim C_L\ge2\dim V=2r$.
\end{proof}

\begin{theorem}[Counterexample to A.2]
\label{C:A2-persistent-arithmetic-counterexample}
There exists an excellent Dedekind domain $R$, of characteristic
zero with $2$ invertible, whose fraction field $K$ is algebraic
over $\mathbf Q$ and satisfies
\[
 \operatorname{cd}_2(K)=1,
 \qquad
 \operatorname{cd}_2(U)=2
 \quad\text{for every nonempty affine open }U\subseteq\operatorname{Spec}R.
\]
In fact, $\operatorname{Pic}(U)/2$ is infinite dimensional over
$\mathbf F_2$ for every nonempty open $U$.
\end{theorem}
\begin{proof}
We start with $F_0=\mathbf Q(\sqrt{-35})$ and
$\mathcal O_{F_0}=\mathbf Z[w]$, where
$w=(1+\sqrt{-35})/2$ satisfies $w^2-w+9=0$.
The discriminant is $-35$. Minkowski's bound
$(2/\pi)\sqrt{35}<4$ shows that the ideal class group is
generated by primes of norm $2$ or $3$; we use
\cite[Section 4, Theorem 4.3, pp.~70--71]{C_MilneANT2020}.
The polynomial of $w$ is irreducible modulo $2$, so $2$ is
inert, its prime ideal is $(2)$, and it has norm $4$.
Modulo $3$ the polynomial has the distinct roots $0$ and $1$.
Thus the primes of norm $3$ are
$P=(3,w)$ and $\overline P=(3,w-1)$. Since $N(w)=9$ and
$w$ is a unit at $\overline P$, we have $(w)=P^2$.
Moreover $P$ is not principal: an element of norm $3$ would
give integers $x,y$ with
\[
 3=N(x+yw)=x^2+xy+9y^2
       =\left(x+\frac y2\right)^2+\frac{35}4y^2,
\]
which is impossible if $y\ne0$ and is impossible for $y=0$
as well. Therefore the class group is $C_2$, generated by
$P$. Inverting $2$ changes nothing because its prime is
principal, so $\dim C_{F_0}=1$.

Enumerate the rational primes, starting with $2$. Having chosen
$F_n$, let $S_n$ consist of all its primes above the first $n+1$
rational primes. Apply Lemma~\ref{C:A2-quadratic-step} to obtain
$F_{n+1}/F_n$, and continue. Set
\[
 K=\bigcup_{n\ge0}F_n,
 \qquad R=\bigcup_{n\ge0}\mathcal O_{F_n}[1/2].
\]
Every fixed rational prime is, from some stage onward, unramified
and inert at every prime above it at every subsequent step.
Consequently it has only finitely many primes above it in $K$,
and all ramification indices along such a prime are bounded.
The ring $R$ is the full integral closure of $\mathbf Z[1/2]$
in $K$. Its localizations at nonzero primes are discrete valuation
rings: they are unions of compatible valuation rings, and bounded
ramification makes the value group discrete. Every nonzero element
has only finitely many prime divisors, by finite splitting.

For completeness, these two facts prove noetherianity. If $I$ is
a nonzero ideal, choose $0\ne a\in I$ and list the finitely many
maximal ideals $P_1,\ldots,P_s$ containing $a$. At each $P_j$
choose $b_j\in I$ attaining the least valuation of a nonzero
element of $I$. Then $J=(a,b_1,\ldots,b_s)$ and $I$ have the
same localizations at all maximal ideals, so $I=J$.
Thus $R$ is Dedekind. This is also the finite-splitting and
discrete-valuation criterion in \cite[Theorem and Corollary,
pp.~40--42]{C_PopescuVraciu1985}. A Dedekind domain with fraction
field of characteristic zero is excellent
\cite[Proposition 15.53.3(4), Tag 07QS]{C_StacksArithmetic}.

We check the generic cohomological dimension with arbitrary
coefficients. Every finite extension $T/K$ has the form $EK$
for a number field $E$ containing some $F_n$. A class in
$\operatorname{Br}(T)[2]$ descends to a Brauer class of order
dividing two over such an $E$, after increasing the finite
stage if necessary. Only finitely many local invariants of
this class are nonzero. At a prime $w$ of $E$ lying above a
prime of $F_n$, the local degrees in $EF_m/E$ have unbounded
two-adic valuation, uniformly over all prolongations of $w$.
To see this precisely, first enlarge $n$ so that the underlying
rational prime is among the prescribed primes. Every local
degree in $F_m/F_n$ at that prime is then $2^{m-n}$. The tower
degree identity in a common completion gives
\[
 v_2\bigl([\,(EF_m)_{w_m}:E_w\,]\bigr)
 \ \ge\ m-n-v_2\bigl([\,E_w:(F_n)_v\,]\bigr).
\]
Thus a single sufficiently large $m$ makes all local restriction
degrees even at the finite support of the Brauer class.
All archimedean places are complex. Restriction multiplies
each local invariant by its local degree, and the
Brauer--Hasse--Noether injection then shows that the class
vanishes over $EF_m$. We have proved
\[
            \operatorname{Br}(T)[2]=0
                 \quad\text{for every finite extension }T/K.
\]
The precise local-global and restriction inputs are
\cite[Chapter III, Proposition 1.8, p.~100; Chapter VII,
Theorem 7.1 and Proposition 7.2, pp.~219--220; Chapter VIII,
Theorem 4.2, p.~241]{C_MilneCFT2020}.

As an algebraic extension of a totally imaginary number field,
$K$ initially has $\operatorname{cd}_2(K)\le2$; we use
\cite[Proposition 8.3.18, with all primes in $S$]{A_NSW}
and closed-subgroup monotonicity. For a finite discrete
$\mathbf F_2[\operatorname{Gal}(K^{\mathrm{sep}}/K)]$-module $M$,
choose a finite Galois extension $T/K$ trivializing it.
A finite sum of regular permutation modules surjects onto $M$.
Its second cohomology vanishes by Shapiro and
$H^2(T,\mathbf F_2)=H^2(T,\mu_2)=\operatorname{Br}(T)[2]=0$.
The kernel has zero third cohomology by the preliminary bound two.
The long exact sequence gives $H^2(K,M)=0$. Filtered colimits
and torsion-coefficient d\'evissage give $\operatorname{cd}_2K\le1$.
There are nonzero primes of $R$, by lying over any nondyadic
prime of $\mathbf Z[1/2]$. A uniformizer in the discrete
valuation ring $R_P$ has odd valuation in $K$, so it is not a
square. Kummer theory gives a nonzero class in $H^1(K,\mu_2)$.
Thus $\operatorname{cd}_2K=1$.

Invertible modules and their isomorphisms descend through
filtered ring colimits, and taking the cokernel of multiplication
by two commutes with such colimits. Hence
\[
 \operatorname{Pic}(R)/2
     =\varinjlim_n \operatorname{Pic}(\mathcal O_{F_n}[1/2])/2.
\]
All transition maps are injective, and the dimensions are at least
$2^n$, by Lemma~\ref{C:A2-quadratic-step}. The resulting vector
space is therefore infinite dimensional. Every nonempty open
of the noetherian integral curve $\operatorname{Spec}R$ removes
only finitely many closed points. Its Picard group is the quotient
of $\operatorname{Pic}(R)$ by the classes of those points.
After reduction modulo two this kills at most a finite-dimensional
subspace. Consequently $\operatorname{Pic}(U)/2$ remains infinite
dimensional for every such $U$.

The Kummer sequence injects $\operatorname{Pic}(U)/2$ into
$H^2(U,\mu_2)$, so every nonempty affine open has
$\operatorname{cd}_2(U)\ge2$. The generic field is nonorderable,
and every closed-point residue field has positive characteristic.
Corollary~\ref{A:generic-plus-dimension},
the generic-field-plus-dimension bound for noetherian schemes,
gives the reverse
inequality $\operatorname{cd}_2(U)\le\operatorname{cd}_2K+1=2$.
This proves both the asserted exact dimensions and the failure
of the full eventual-equality quantifiers of A.2.
\end{proof}

\paragraph{Finite verification and scope.}
The accompanying exact-arithmetic program
\nolinkurl{parts/C/code/check_quadratic_tower.py} checks the wreath-group
identities for $\dim V=1,2,3$, the transfer from its coset definition,
every local Frobenius-square implication, the inertia complements,
and the initial norm calculation and reduced forms. Its saved output
\nolinkurl{parts/C/output/quadratic_tower_certificate.json} records
$111{,}784$ passed assertions in $0.068522113$ seconds. These finite checks
supplement the proofs above; the infinite tower is justified by the
quadratic-step lemma, Chebotarev, and the explicit noetherianity and
Brauer arguments. No positive Iwasawa $\mu$-invariant or finite-stage
class-rank growth without persistence is assumed.
